\documentclass[exercice,a4]{thedocument}%
\usepackage{texmaths-tikz}%
\usepackage{pstricks}%
\usepackage{pst-bezier}%
\usepackage{pstricks-add}%
\usepackage{pst-eucl}%
\usepackage{pst-func}%
\usepackage{pst-tree}%
\usepackage{pst-circ}%
\usepackage{tikz}%
\usepackage{pgfplots}%
\usepackage{tkz-tab}%
\usepackage{tikzpagenodes}%
\usepackage{tabularray}%
\usepackage{tabularx}%
\usepackage{multirow}%
\usepackage{multicol}%
\usepackage{scratch3}%
\usepackage{enumitem}%
\usepackage[np]{numprint}%
\usepackage{pas-tableur}%
\usepackage{listings}%
\usepackage{graphicx}%
\usepackage{adjustbox}%
\usepackage{needspace}%
\usepackage{xcolor}%
\usepackage{eurosym}%
\usepackage{diagbox}%
\usepackage{colortbl}%
\usepackage{multido}%
\usepackage{icomma}%
\usepackage[french]{babel}%
\graphicspath{{../brevet/2026/}}%
\startdatakeys
    \source{26GENMATAA1}
    \BO{2026}
    \cycle{4}
    \competencesDuSocle{CA,RA,CO}
    \objectifsApprentissage{C4A1aa,C4A2ad,C4B2cb,C4D2aa,C4D3ac}
    \motsCles{programme de calcul,tableur,expression littérale,équation produit}
\enddatakeys
\definecolor{gris}{gray}{0.5}%
\providecommand{\wideparen}[1]{\overset{\frown}{#1}}%
\begin{document}%
\providecommand{\degres}{^\circ}%
On considère le \textbf{programme A} suivant :

\begin{center}
	\begin{tikzpicture}[
		box/.style={draw, rounded corners, align=center, minimum width=3.6cm, minimum height=0.9cm},
		>={Stealth[length=2.5mm]}
		]
		\node[box] (start) at (3,6) {Choisir un nombre};
		\node[box] (carre) at (0,4.5) {Calculer son carré};
		\node[box] (mult7) at (6,4.5) {Multiplier par 7};
		\node[box] (mult3) at (0,3) {Multiplier par 3};
		\node[box] (add) at (3,1.5) {Additionner les deux nombres};
		\node[box] (sous) at (3,0) {Soustraire 6};

		\draw[->] (start) -- (carre);
		\draw[->] (start) -- (mult7);
		\draw[->] (carre) -- (mult3);
		\draw[->] (mult3) -- (add);
		\draw[->] (mult7) -- (add);
		\draw[->] (add) -- (sous);
	\end{tikzpicture}
\end{center}

\begin{minipage}[t]{0.6\linewidth}
	\begin{enumerate}
	\item Appliquer le \textbf{programme A} au nombre 5.

	\item On utilise un tableur pour trouver les résultats correspondants à quelques nombres comme l'indique le tableau ci-contre.

	\medskip

	Parmi les quatre formules ci-dessous, recopier celle qui a été saisie dans la cellule \textsf{B2}, puis étirée vers le bas afin de calculer les résultats donnés par le programme A.

	\medskip

	\emph{Aucune justification n'est attendue.}

	\begin{tblr}{colspec={XX},hlines, vlines,cells={cmd={\sffamily \small}}}
			= 3 * A2 * 2 + 7 * A2 - 6 & = 3 * 1 * 1 + 7 * 1 - 6 \\
			= 3 * A2 * A2 + 7 * A2 - 6& = 3 * A2 * 2 - 7 * A2 + 6 \\
		\end{tblr}
\end{enumerate}
\end{minipage}
\hfill
\begin{tblr}[t]{hlines={darkgray}, vlines, columns={c, cmd={\sffamily \small}},
			column{2,3}={wd=2.7cm,colsep=2pt},
			row{1}={rowsep=0pt,ht=7mm, bg=gray9},
			row{2-Z}={rowsep=1pt},
			column{1}={colsep=0pt,7mm,bg=gray9},
			cell{1}{1}={gray9},
			cell{2}{2-3}={cmd={\sffamily \bfseries}}}
			\SetCell{halign=r,valign=b}\tikz[baseline={(0,0.15)}]{\clip(0,0) rectangle (0.6,0.5) ;\fill[gray!50] (0,0)--(0.5,0.5)--(0.5,0);}
			   & A                & B                       \\
			1  & Nombre de départ & Résultat du programme A \\
			2  & -3,5             & 6,25                    \\
			3  & -3               & 0                       \\
			4  & -2,5             & -4,75                   \\
			5  & -2               & -8                      \\
			6  & -1,5             & -9,75                   \\
			7  & -1               & -10                     \\
			8  & -0,5             & -8,75                   \\
			9  & 0                & -6                      \\
			10 & 0,5              & -1,75                   \\
			11 & 1                & 4                       \\
			12 & 1,5              & 11,25                   \\
			13 & 2                & 20                      \\
	\end{tblr}
\begin{enumerate}[start=3]
	\item \`A l'aide du tableur, donner une valeur pour laquelle le \textbf{programme A} donne 0.

	\textit{Aucune justification n'est attendue.}

	\item Si on note $x$ le nombre de départ, donner une expression littérale du programme A en fonction de~$x$.
\end{enumerate}

\newpage

On considère maintenant le \textbf{programme B} suivant :

\begin{center}
	\begin{tikzpicture}[
		box/.style={draw, rounded corners, align=center, minimum width=3.6cm, minimum height=0.9cm},
		>={Stealth[length=2.5mm]}
		]
		\node[box] (start) at (3,5) {Choisir un nombre};
		\node[box] (mult3) at (0,3.5) {Multiplier par 3};
		\node[box] (add3) at (6,2.75) {Additionner 3};
		\node[box] (sous2) at (0,2) {Soustraire 2};
		\node[box] (multdeux) at (3,0.5) {Multiplier les deux nombres};

		\draw[->] (start) -- (mult3);
		\draw[->] (start) -- (add3);
		\draw[->] (mult3) -- (sous2);
		\draw[->] (sous2) -- (multdeux);
		\draw[->] (add3) -- (multdeux);
	\end{tikzpicture}
\end{center}

\begin{enumerate}[resume]
	\item Appliquer le \textbf{programme B} au nombre 5.

	\item Si on note $x$ le nombre de départ, donner une expression littérale du \textbf{programme B} en fonction de~$x$.

	\item Mathis affirme que, quel que soit le nombre qu'il choisit, il trouvera le même résultat avec le \textbf{programme A} et le \textbf{programme B}. A-t-il raison ? Justifier.

	\item Résoudre l'équation $(3x-2)(x+3) = 0$.

	En déduire les valeurs de $x$ pour lesquelles les programmes A et B donnent 0.
\end{enumerate}

\bigskip
\medskip
\showthesource
\end{document}
